\chapter{Parent-origin минимального бинарного incidence-носителя}
\label{ch:v8-up-down-binary-carrier-parent-origin}

\section{Статический центр уже существует}

В расширенном endpoint-носителе оба цветных типа уже представлены
центральными проекторами. Поэтому классическая бинарная алгебра
\begin{equation}
 \mathcal A_{ud}^{\rm cl}=\mathbb CP_u\oplus\mathbb CP_d,
 \qquad P_uP_d=0,\qquad \dim\mathcal A_{ud}^{\rm cl}=2
 \label{eq:v8-binary-carrier-existing-center}
\end{equation}
не требует ещё одной абстрактной копии $\mathbb C^2$. Это положительный
kinematic origin только для двух меток. Он не содержит переключателя между
ними.

На бинарном quotient гиперзаряд действует как
\begin{equation}
 Y_{ud}=\begin{pmatrix}5/3&0\\0&2/3\end{pmatrix}.
 \label{eq:v8-binary-carrier-hypercharge}
\end{equation}
Его инвариантный коммутант диагонален:
\begin{equation}
 [Y_{ud},X]=0
 \quad\Longleftrightarrow\quad
 X=\begin{pmatrix}x_u&0\\0&x_d\end{pmatrix},
 \qquad \dim\{Y_{ud}\}'=2.
 \label{eq:v8-binary-carrier-invariant-commutant}
\end{equation}
Следовательно, invariant Hamiltonian старого endpoint-блока не содержит
$\sigma_x$ или $\sigma_y$ и не меняет бинарную популяцию.

\section{Ковариантный скачок не является invariant-стрелкой}

Операторы перехода
\begin{equation}
 L_\downarrow=|d\rangle\langle u|,
 \qquad L_\uparrow=L_\downarrow^\dagger
 \label{eq:v8-binary-carrier-transition-operators}
\end{equation}
имеют определённые $U(1)$-веса:
\begin{equation}
 [Y_{ud},L_\downarrow]=-L_\downarrow,
 \qquad [Y_{ud},L_\uparrow]=+L_\uparrow.
 \label{eq:v8-binary-carrier-transition-charges}
\end{equation}
Поэтому invariant Hom между endpoint-типами нулевой, но GKSL-диссипатор
одного charged jump всё же gauge-ковариантен. При
$U_\theta L_\downarrow U_\theta^\dagger=e^{-i\theta}L_\downarrow$
фаза сокращается:
\begin{equation}
 \mathcal D_{e^{-i\theta}L_\downarrow}(X)
 =\mathcal D_{L_\downarrow}(X).
 \label{eq:v8-binary-carrier-dissipator-covariance}
\end{equation}
Это различие существенно: симметрия допускает условный открытый процесс,
но не создаёт его микроскопический оператор.

\section{Старый 42-frame не действует на новый switch}

Полный старый шумовой кадр удовлетворяет
\begin{equation}
 \mathcal F_{42}\subset\operatorname{End}(H_{21}).
 \label{eq:v8-binary-carrier-old-frame-ambient}
\end{equation}
После расширения новым бинарным off-diagonal сектором его каноническое
вложение остаётся нулевым на этом секторе:
\begin{equation}
 \iota(F_a)=F_a\oplus0_{\rm bit},
 \qquad P_{\rm bit}\iota(F_a)=0=\iota(F_a)P_{\rm bit}.
 \label{eq:v8-binary-carrier-old-frame-zero-extension}
\end{equation}
Ни линейная оболочка, ни trace-dual whitening не создают отсутствующий
$L_\downarrow$. Следовательно, прежняя 42-скачковая примитивность не
переносится автоматически на новый controlled incidence.

\section{Минимальный компенсирующий environment-тип}

Gauge-инвариантный системно-средовой гамильтониан должен компенсировать
заряд $-1$ оператора $L_\downarrow$. Условная минимальная форма есть
\begin{equation}
 H_{\rm int}^{ud}
 =g\left(L_\downarrow\otimes b_+
 +L_\uparrow\otimes b_+^\dagger\right),
 \label{eq:v8-binary-carrier-compensated-interaction}
\end{equation}
где
\begin{equation}
 [Y_{\rm env},b_+]=+b_+,\qquad
 [Y_{\rm tot},H_{\rm int}^{ud}]=0.
 \label{eq:v8-binary-carrier-charge-compensation}
\end{equation}
Таким образом, нужен environment-transition заряда $+1$, например переход
из нейтрального вакуума в Real-сопряжённую charged-singlet линию.
Ковариантные Стайнспринг-дилатации именно так переносят представление jump-
пространства на среду \cite{v8-binary-carrier-haapasalo}.

В старом $H_{21}$ charged-singlet тип $Y=-1$ имеет кратность три; его
Real-сопряжение предоставляет три кандидата заряда $+1$:
\begin{equation}
 \dim\mathcal K_{+1}^{\rm old}=3.
 \label{eq:v8-binary-carrier-charged-multiplicity}
\end{equation}
Но это семейный триплет. Для стандартных генераторов $J_1,J_2,J_3$ имеем
\begin{equation}
 \bigcap_{a=1}^{3}\ker J_a=\{0\}.
 \label{eq:v8-binary-carrier-family-fixed-zero}
\end{equation}
Следовательно, старый charged carrier не содержит канонической family-
singlet линии $b_+$. Выбор одного из трёх состояний ломает семейную
ковариантность, а усреднение плотности не создаёт линейный interaction-
оператор.

\section{Энергия и состояние среды}

Даже после добавления подходящей charged-singlet линии энергосохранение
требует резонанса
\begin{equation}
 E_{+1}-E_0=E_u-E_d=7\gamma.
 \label{eq:v8-binary-carrier-resonance}
\end{equation}
Для downward-предела среда должна начинаться в нейтральном вакууме и не
возвращать возбуждение:
\begin{equation}
 \rho_{\rm env}=|0\rangle\langle0|,
 \qquad b_+^\dagger|0\rangle=0
 \quad\text{в принятой ориентации annihilation/creation}.
 \label{eq:v8-binary-carrier-environment-state}
\end{equation}
Ориентация оператора может быть переобозначена, но физически обязательны
низкотемпературное состояние, резонанс и ненулевая spectral density. Ни
одно из них не следует из статического endpoint-центра.

Минимальный недостающий пакет равен
\begin{equation}
 \mathcal D_{\rm miss}=
 \{L_\downarrow,\mathcal K_{+1}^{\rm sing},
 H_{\rm int}^{ud},\rho_{\rm env},
 E_{+1}-E_0=7\gamma,\kappa_{ud}\}.
 \label{eq:v8-binary-carrier-missing-data}
\end{equation}

Итоговые реестры:
\begin{equation}
 N_{\rm static\ center}=4/4,\qquad
 N_{\rm covariant\ transition}=6/6,\qquad
 N_{\rm microscopic\ parent}=0/6.
 \label{eq:v8-binary-carrier-ledgers}
\end{equation}

\begin{theorem}[Статический carrier без микроскопического switch]
Проекторы $P_u,P_d$ уже задают двумерную классическую бинарную алгебру.
Оператор $L_\downarrow$ имеет гиперзаряд $-1$, и его GKSL-диссипатор
gauge-ковариантен. Однако invariant endpoint-Hom нулевой, старый 42-frame
аннулирует новый off-diagonal сектор, а доступный charged-singlet mediator
является семейным триплетом без неподвижной линии. Поэтому микроскопический
parent бинарной релаксации не получен.
\end{theorem}

\begin{nogocriterion}[Ковариантность dissipator не создаёт dilation]
Нельзя из фазовой ковариантности $\mathcal D_{L_\downarrow}$ заключать, что
$L_\downarrow$ содержится в старом шумовом кадре. Нельзя выбирать одну из
трёх charged-singlet environment-линий без family-селектора. Даже новый
mediator не закрывает динамику без резонансной щели, состояния среды и
физической скорости.
\end{nogocriterion}

\begin{protocol}\begin{enumerate}
\item статическая бинарная центральная алгебра: \textbf{уже существует};
\item её размерность: \textbf{$2$};
\item invariant off-diagonal Hom: \textbf{$0$};
\item заряд $L_\downarrow$: \textbf{$-1$};
\item его GKSL-диссипатор gauge-ковариантен: \textbf{да};
\item старый 42-frame содержит $L_\downarrow$: \textbf{нет};
\item требуемый environment-заряд: \textbf{$+1$};
\item старых charged-singlet кандидатов: \textbf{$3$};
\item family-invariant mediator line: \textbf{$0$};
\item требуемая резонансная щель: \textbf{$7\gamma$};
\item состояние и скорость среды выведены: \textbf{нет};
\item static center: \textbf{$4/4$};
\item covariant transition shape: \textbf{$6/6$};
\item microscopic parent-origin: \textbf{$0/6$};
\item следующий гейт: \textbf{архитектура charged environment mediator}.
\end{enumerate}\end{protocol}

Машинный сертификат сохранён в
\path{s2t_v8_baryon_c0_singlet_triplet_central_gap_isotypic_channel_extra_edge_mass_two_source_cubic_incidence_up_down_binary_selector_minimal_typed_discriminator_binary_carrier_parent_origin_gate_results.json}.

\begin{thebibliography}{9}
\bibitem{v8-binary-carrier-haapasalo} E.~Haapasalo,
\emph{The Choi--Jamio\l{}kowski Isomorphism and Covariant Quantum Channels},
Quantum Stud. Math. Found. 8 (2021) 1; arXiv:1906.11442.
\bibitem{v8-binary-carrier-lindblad} G.~Lindblad,
\emph{On the Generators of Quantum Dynamical Semigroups},
Commun. Math. Phys. 48 (1976) 119.
\end{thebibliography}